\section{Extremal weights from twisted differential forms}
\label{V:sec-extremal}

The transfer theorem becomes equivariant when its arc is a generic
one-parameter orbit.  Its bounded input order comes from the fixed
shift between a cohomological weight and the twist of a differential
form.  The output sections approximate every exposed vertex of the
anticanonical weight polytope, after which multiplication clears the
weights.

The cohomological assumptions here are $q(W)=0$ and
$\chi(W,\mathcal O_W)\ne0$.  The virtual character determines the
source coefficient directly from this Euler characteristic.  The
rest of the argument follows the source-weight construction of
Section~\ref{subsec:valuation-extremal-weights}, using pseudoeffectivity of the
inverse determinant and approximate arc transfer to approach the
vertices.  Thus the proof does not require vanishing of all the higher
structure-sheaf cohomology groups or exact attainment of a source weight.

\begin{proposition}\label{V:extremal-invariants}
Let $W$ be smooth connected projective with $q(W)=0$ and
$\chi(W,\mathcal O_W)\ne0$.  Suppose $A=-K_W$ has a smooth
semipositive Hermitian metric.  Let an algebraic torus $T$ act
faithfully on $W$, and suppose an equivariantly linearized very ample
line bundle is available.  For the natural linearization on $A$,
there is an integer $r>0$ with $H^0(W,rA)^T\ne0$.
\end{proposition}

\begin{proof}
The point case is immediate.  In the real character space of $T$ put
\[
 \mathcal P=\operatorname{conv}
 \{\hbox{fiber characters of $A$ on the components of $W^T$}\}.
\]
The fixed locus is nonempty and has finitely many components, so this
is a rational polytope.  Fix an integral cocharacter
$\lambda:\mathbb C^*\to T$.  For a nontrivial action let $S_+$ be the
source component, whose attracting set for parameter tending to zero
is dense, and let $\mu_+$ be the $\lambda$-weight of $A$ there.
The Bia\l ynicki-Birula decomposition \cite{BB1973} shows that all
normal weights at $S_+$ are positive and that every other fixed
component has a negative normal direction.  Following such a direction
backwards along an orbit moves to a component of strictly larger
weight for the fixed very ample linearization.  For any nef
linearized line bundle, its weight at the zero end minus its weight
at the infinity end is the degree on the parametrized orbit
$\mathbb P^1$ and is nonnegative.  Iterating the backwards move
therefore terminates at the source and gives
\begin{equation}\label{V:source-support}
 \mu_+=\max_{\alpha\in\mathcal P}\langle\alpha,\lambda\rangle.
\end{equation}
Here $S_+$ is $T$-stable by uniqueness and contains a $T$-fixed point.
Faithfulness makes the $\lambda$-action nontrivial for every nonzero
cocharacter.  Since the natural anticanonical weight is the sum of
tangent weights, $\mu_+>0$ for every such cocharacter.  Rational
polyhedral separation implies that $0$ is in the interior of
$\mathcal P$.  If $T$ is trivial, use the zero cocharacter, $S_+=W$,
and $\mu_+=0$ in the argument below.

\subsection*{The largest cohomological weight}

Apply the equivariant holomorphic fixed-point formula
\cite{AtiyahSegal1968,AtiyahSinger1968} to the virtual character
$\sum_j(-1)^j H^j(W,mA)$.  A fixed component $U$ contributes
\[
 t^{m\mu_U}\int_U e^{m c_1(A)|_U}\operatorname{td}(T_U)
       \prod_\nu(1-t^{-d_\nu}e^{-x_\nu})^{-1},
\]
where $d_\nu\ne0$ are its normal weights and $x_\nu$ the corresponding
Chern roots.  Expand at infinity.  A positive $d_\nu$ contributes a
series with highest exponent zero, whereas a negative $d_\nu$
contributes only strictly negative exponents.  Thus for $m\ge0$ the
coefficient of $t^{m\mu_+}$ is exactly
\[
 Q(m)=\chi(S_+,mA|_{S_+}).
\]
Indeed every other component has a negative normal direction and
$\mu_U\le\mu_+$.  At $m=0$, the connected torus acts trivially on
$H^j(W,\mathcal O_W)$: it acts trivially on ordinary cohomology by
homotopy, hence on its Hodge summands.  The constant coefficient is
therefore $\chi(W,\mathcal O_W)$.  It follows that
$Q(0)=\chi(W,\mathcal O_W)\ne0$.
Riemann--Roch makes $Q$ a polynomial, so it is nonzero for all
sufficiently large integers $m$.

For some fixed $j$, infinitely many $H^j(W,mA)$ therefore contain
weight $m\mu_+$.  Average a K\"ahler form over the compact circle of
$\lambda$.  Semipositive hard Lefschetz
\cite{DPS2001}, applied to $(m+1)A$, gives the equivariant surjection
\[
 H^0(W,\Omega_W^{n-j}\otimes(m+1)A)
 \longrightarrow H^j(W,K_W+(m+1)A)=H^j(W,mA),
 \qquad n=\dim W.
\]
The equality uses the natural anticanonical linearization; its shift
is $m+1$.  Projecting to the desired circle weight gives a nonzero
holomorphic eigenform of weight $m\mu_+$ for each of infinitely
many $m$.

\subsection*{Saturated determinants and the fixed order}

We use the determinant method of Lazi\'c--Peternell
\cite[Lemma~4.1]{LP2018} and its anticanonical adaptation in
\cite[Lemma~5.1]{LMPTX2023}.  Those published statements assume
that the twisting line is numerically nontrivial.  Here the twists
are already positive and unbounded, so the elementary construction
below also handles the numerically trivial case.  The weights of
its wedges, including the rank shift, are computed locally.

Put $p=n-j$ and take the generic span of these forms as lines in
$\Omega_W^p$ over $\mathbb C(W)$.  Let its rank be $s>0$.  Wedge
$s$ generic basis forms.  The saturation of the resulting generic
line in $\bigwedge^s\Omega_W^p$ is a line bundle $M$, with its induced
$\lambda$-linearization.  More concretely, the wedge has some twist
$a_1A$; removing its common divisorial zero $D_1$ gives
$M=D_1-a_1A$ and a map into $\bigwedge^s\Omega_W^p$ nonvanishing in
codimension one.  The saturation is reflexive of rank one, hence
invertible on the smooth variety.  Its map extends across codimension
two, and regular generic-line sections factor through it because
their coefficients have no poles in codimension one.

The negative divisor $-M$ is pseudoeffective.  When $p>0$, the
exterior/tensor inclusion places this saturated line in a positive
tensor power of $\Omega_W^1$.  The theorem of
Lazi\'c--Matsumura--Peternell--Tsakanikas--Xie
\cite[Theorem~4.1]{LMPTX2023}, following Ou's movable-slope theorem
\cite[Theorem~1.4]{Ou2023}, applies to the smooth zero-boundary pair
with nef anticanonical divisor.  If $p=0$, the line is the underlying
trivial bundle, so the same conclusion holds directly.

Complete each eigenform of increasingly large twist to a generic
basis using $s-1$ members of the fixed finite basis, then take an
infinite subsequence with the same choice of these members.  This
gives nonzero eigensections and effective divisors
\begin{equation}\label{V:weighted-determinants}
 N_i\sim M+m_iA,
 \qquad\hbox{section weight }(m_i-s)\mu_+,
 \qquad 0<m_1<m_2<\cdots.
\end{equation}
The exponent $m_i$ is the sum of the $s$ form twists $m+1$; the
subtracted $s$ in the weight is therefore essential.

Let $v$ be the order along the generic orbit arc
$t\mapsto\lambda(t)\xi$, where $\xi$ is the generic point regarded
as a $\mathbb C(W)$-valued point.  Projectivity extends it to
$\operatorname{Spec}\mathbb C(W)[[t]]$, and its fraction point still
maps to the generic point.  Its center is in $S_+$.
Lemma~\ref{V:order-weight} applies with the declared left action on
sections.  The order is the fiber weight minus the section weight.

Writing $\mu_M$ for the fiber weight of $M$ on the source, equation~\eqref{V:weighted-determinants} and Lemma~\ref{V:order-weight} give
\[
 v(N_i)=\mu_M+m_i\mu_+-(m_i-s)\mu_+
       =\mu_M+s\mu_+.
\]
These are bounded unnormalized orders.  Theorem~\ref{V:approximate-arc}
therefore supplies sections of $rA$ with $v/r$ as small as desired.

\subsection*{Approximating vertices and multiplying characters}

Decompose each such section into $T$-weight vectors.  At least one
nonzero component has order no greater than the original section,
since the order of a sum is at least the minimum of its summands'
orders.  If its character is $\alpha$, Lemma~\ref{V:order-weight} gives
\[
 \frac{\langle\alpha,\lambda\rangle}{r}
     =\mu_+-\frac{v}{r}.
\]
Every normalized character $\alpha/r$ of an actual section belongs
to $\mathcal P$: applying nonnegativity of order to each integral
cocharacter gives all the rational supporting-half-space inequalities
from \eqref{V:source-support}.

Choose one integral exposing direction for each vertex of
$\mathcal P$.  In such a direction the maximum is unique.  Compactness
of the polytope then turns approach to the maximal value into approach
to that vertex.  Thus every vertex is a limit of normalized characters
of actual sections.  Since zero is interior, choose finitely many
of these normalized characters close enough to the vertices that
zero remains in their convex hull.  Lemma~\ref{V:convex-product} now gives a nonzero invariant section
of a positive multiple of $A$.
For the trivial torus the transfer theorem has already given the
required section.  This proves the proposition.
\end{proof}
